\section{Вариационные методы решения краевых задач и определения собственных значений}
\label{sec:variational-bvp-eigen}

\subsection{Слабая постановка краевой задачи}

Рассмотрим модельную эллиптическую задачу
\[
-\nabla\cdot(k\nabla u)=f\quad\text{в }\Omega,
\qquad
u=0\quad\text{на }\Gamma_D,
\qquad
k\frac{\partial u}{\partial n}=g\quad\text{на }\Gamma_N.
\]
Умножая уравнение на тестовую функцию и интегрируя по частям, получаем: найти $u\in V$,
\[
V=\{v\in H^1(\Omega):v|_{\Gamma_D}=0\},
\]
такую, что
\[
\boxed{a(u,v)=l(v)\qquad\forall v\in V},
\]
где
\[
a(u,v)=\int_\Omega k\nabla u\cdot\nabla v\,dx,
\qquad
l(v)=\int_\Omega fv\,dx+\int_{\Gamma_N}gv\,ds.
\]
Слабая постановка требует меньшей гладкости решения и является основой методов Ритца, Галёркина и МКЭ.

\textbf{Теорема Лакса--Мильграма.} Пусть $V$ --- гильбертово пространство, билинейная форма непрерывна
\[
|a(u,v)|\le M\|u\|_V\|v\|_V
\]
и коэрцитивна
\[
a(v,v)\ge\alpha\|v\|_V^2,
\qquad \alpha>0,
\]
а $l\in V^*$. Тогда существует единственное $u\in V$, удовлетворяющее $a(u,v)=l(v)$ для всех $v\in V$.

\subsection{Энергетический функционал и метод Ритца}

Если $a$ симметрична, слабая задача эквивалентна минимизации
\[
J(v)=\frac12a(v,v)-l(v).
\]
Условие $\delta J(u;w)=0$ даёт слабое уравнение. В механике $J$ часто является полной потенциальной энергией.

В \textbf{методе Ритца} выбирают
\[
V_N=\operatorname{span}\{\varphi_1,\ldots,\varphi_N\}\subset V,
\qquad
u_N=\sum_{j=1}^{N}c_j\varphi_j
\]
и минимизируют $J(u_N)$. Получается система
\[
Kc=f,
\qquad
K_{ij}=a(\varphi_j,\varphi_i),
\quad
f_i=l(\varphi_i).
\]

\subsection{Метод Галёркина и оценка ошибки}

В \textbf{методе Галёркина} требуется
\[
a(u_N,v_N)=l(v_N)
\qquad \forall v_N\in V_N.
\]
Для симметричной потенциальной задачи Ритц и Галёркин дают одну систему: Ритц исходит из минимума функционала, Галёркин --- из ортогональности невязки.

\textbf{Лемма Сеа.} При условиях Лакса--Мильграма для галёркинского приближения
\[
\boxed{
\|u-u_N\|_V
\le
\frac{M}{\alpha}
\inf_{v_N\in V_N}\|u-v_N\|_V.
}
\]
То есть метод квазиоптимален: его ошибка не более чем в постоянное число раз превосходит наилучшую ошибку аппроксимации выбранным пространством.

\subsection{Спектральная задача и отношение Рэлея}

Для самосопряжённой спектральной задачи слабая форма имеет вид
\[
a(u,v)=\lambda b(u,v)
\qquad \forall v\in V,
\qquad u\ne0,
\]
где $a$ и $b$ симметричны, $b(u,u)>0$ для $u\ne0$. В механике после дискретизации типична задача
\[
Kx=\omega^2Mx.
\]

Для $v\ne0$ вводится \textbf{отношение Рэлея}
\[
R(v)=\frac{a(v,v)}{b(v,v)}.
\]
Для собственной функции $R(u)=\lambda$.

При стандартных предположениях для положительной самосопряжённой задачи первое собственное значение характеризуется принципом Рэлея:
\[
\boxed{\lambda_1=\min_{v\ne0}R(v)}.
\]
Более общий \textbf{принцип Куранта--Фишера (min--max)} имеет вид
\[
\boxed{
\lambda_k=
\min_{\substack{S\subset V\\ \dim S=k}}
\;\max_{0\ne v\in S}R(v).
}
\]
Эквивалентно, при ортонормированном наборе собственных функций первые значения можно получать последовательной минимизацией при $b$-ортогональности к ранее найденным модам.

\subsection{Метод Рэлея--Ритца}

В конечномерном пространстве $V_N$ задача принимает вид
\[
Kc=\lambda Mc,
\]
где
\[
K_{ij}=a(\varphi_j,\varphi_i),
\qquad
M_{ij}=b(\varphi_j,\varphi_i).
\]
Ненулевое решение существует при
\[
\det(K-\lambda M)=0.
\]

Для согласованных подпространств вариационный принцип даёт важное свойство: ритцевские собственные значения приближают точные \textbf{сверху},
\[
\lambda_k\le \lambda_k^{(N)},
\]
а при вложенных пространствах $V_N\subset V_{N+1}$ обычно
\[
\lambda_k^{(N+1)}\le \lambda_k^{(N)}.
\]
При плотном обогащении пространств и стандартных условиях спектральной теории эти приближения сходятся к истинным собственным значениям.

\subsection{Что важно сказать на экзамене}

Нужно показать цепочку
\[
\boxed{
\text{краевая задача}
\rightarrow
\text{слабая форма}
\rightarrow
\text{Лакс--Мильграм}
\rightarrow
\text{Ритц/Галёркин}
\rightarrow
\text{Сеа}}
\]
и для спектральной части записать
\[
a(u,v)=\lambda b(u,v),
\qquad
R(u)=\frac{a(u,u)}{b(u,u)},
\qquad
Kc=\lambda Mc.
\]
Центральные вариационные результаты --- минимум Рэлея и принцип min--max.

\newpage
